\section{Forward expansions}
\label{sec:forward}

Before inverting anything we need the forward function in the normal form $ax^p(1+\sum_ix^{\delta_i}B_i(\log x))$, and the second question \cite{q2} is a forward problem in its own right. This section shows that a large class of expressions has convergent power--log expansions, and describes the precision-tracking evaluator that computes them with a rigorous remainder class.

\subsection{Convergent power--log expansions}

\begin{definition}
\label{def:pl-expansion}
A function $F$ on $(0,w_0)$ has a \emph{convergent power--log expansion} with support in $\beta_0+S$ ($S$ a finitely generated positive semigroup with gaps $\delta_1,\dots,\delta_m$) if there are polynomials $C_\sigma\in\R[L]$, constants $M,K\ge1$, and integers $d_0,d\ge0$ such that
\begin{equation}
 F(w)=\sum_{\sigma\in S}w^{\beta_0+\sigma}C_\sigma(L)\quad\text{for }0<w<w_1,
 \qquad \|C_\sigma\|_1\le M K^{n(\sigma)},\quad
 \deg C_\sigma\le d_0+d\,n(\sigma),
 \label{eq:pl-expansion}
\end{equation}
where $n(\sigma)$ is the least $|\bk|$ with $\bk\cdot\bdelta=\sigma$, $\|P\|_1$ is the sum of the absolute values of the coefficients of $P$, and the series converges absolutely. Thus $|C_\sigma(L)|\le M K^{n(\sigma)}\M(w)^{d_0+dn(\sigma)}$. Zero coefficients satisfy the degree bound by convention. We write $F\in\mathcal E$. The fixed allowance $d_0$ is essential: it admits a polynomial in $\log w$ at the leading exponent, including $\log w$ itself.
\end{definition}

The coefficient bound makes the tail after any cutoff controllable (Lemma~\ref{lem:tail} below), and it is stable under all the operations we need.

\begin{lemma}[Tail of a convergent expansion]
\label{lem:tail}
Let $F\in\mathcal E$ as in \eqref{eq:pl-expansion}, $\delta_*=\min_i\delta_i$, and $H>0$. Then
\begin{equation}
 F(w)-\sum_{\sigma<H}w^{\beta_0+\sigma}C_\sigma(L)=\OO\bigl(w^{\beta_0+H}\M(w)^{d_0+Nd}\bigr),\qquad N=\lceil H/\delta_*\rceil.
 \label{eq:tail}
\end{equation}
More precisely, the omitted part equals $w^{\beta_0+\sigma_*}C_{\sigma_*}(L)+\oo(w^{\beta_0+\sigma_*})$ where $\sigma_*=\min\{\sigma\in S:\sigma\ge H\}$.
\end{lemma}
\begin{proof}
Split the omitted terms into those with $n(\sigma)\le N-1$ and the rest. The first group is finite (Lemma~\ref{lem:finite}) and each term is $\OO(w^{\beta_0+H}\M^{d_0+(N-1)d})$. For the rest, the number of weights whose shortest representation has length $n$ is at most $m^n$, and $\sigma\ge n\delta_*$, so their sum is bounded by
\[
 Mw^{\beta_0}\M^{d_0}\sum_{n\ge N}\bigl(mKw^{\delta_*}\M^d\bigr)^n
 =\OO(w^{\beta_0+N\delta_*}\M^{d_0+Nd}).
\]
Since $N\delta_*\ge H$, this proves the bound. For the refined statement let $\sigma_{**}$ be the next element of $S$ after $\sigma_*$, which exists by local finiteness when $S$ has a generator, and choose $\sigma_*<H'<\sigma_{**}$. Applying the bound at $H'$ shows that everything beyond $\sigma_*$ is $\OO(w^{\beta_0+H'}\M^{d_0+N'd})=\oo(w^{\beta_0+\sigma_*})$. The difference $\sigma_{**}-\sigma_*$ need not be at least $\delta_*$. Finite support with no further terms has zero tail.
\end{proof}

\begin{theorem}[Closure]
\label{thm:closure}
The class $\mathcal E$ contains the constants and $w$ and is closed under sums, products, and the following operations, with the supports and gaps indicated.
\begin{enumerate}[leftmargin=1.8em]
\item If $F\in\mathcal E$ has leading block $cw^\beta$ with a \emph{constant} $c>0$, then $F^a\in\mathcal E$ for every real $a$ and $\log F\in\mathcal E$; their supports are contained in $a\beta+S'$ and $\{0\}\cup S'$ respectively, with block $\log c+\beta L$ at exponent $0$ in the logarithm. Here $S'$ is a finitely generated positive semigroup containing the exponents of $F/(cw^\beta)-1$; the proof establishes that such a semigroup exists.
\item If $F\in\mathcal E$ has no negative-exponent blocks and its exponent-$0$ block is $c+kL$ with $k\in\R$, then $e^F\in\mathcal E$ with leading block $e^cw^k$. The argument $F$ may grow logarithmically.
\item If $F\in\mathcal E$ has $\beta_0\ge0$ and constant exponent-$0$ block $c$, and $\phi$ is analytic at $c$, then $\phi\circ F\in\mathcal E$.
\end{enumerate}
\end{theorem}
\begin{proof}
First note that a positive-valuation expansion $U$ can be represented with support in a positive semigroup based at $0$: if necessary, add its positive leading exponent as a generator and enlarge the constants in \eqref{eq:pl-expansion}. Choose one shortest multi-index $\bk$ for each resulting weight $\sigma$, put $D=d_0+d$, and introduce
\[
 T_{ij}=w^{\delta_i}L^j,\qquad 1\le i\le m,\quad 0\le j\le D.
\]
For $n=|\bk|\ge1$, every $w^\sigma L^s$ with $s\le d_0+dn\le Dn$ is a product of $n$ such parameters: distribute the $s$ logarithmic factors among the $n$ generator factors, each receiving at most $D$. Fix one such distribution for every monomial. The coefficient-norm bound makes the resulting series analytic in a sufficiently small polydisc, since its terms of degree $n$ have total absolute coefficient at most $M(mK)^n$. Analytic composition with $(1+U)^a$, $\log(1+U)$, $e^U$, or $\phi(c+U)$ is therefore legitimate. Cauchy estimates give exponential coefficient bounds in the new parameters, and their logarithmic degree grows at most linearly in parameter degree.

Regrouping by weight preserves bounds of the form \eqref{eq:pl-expansion}: a representation of weight $\sigma$ has length at most $\sigma/\delta_*$, whereas its shortest length is at least $\sigma/\delta_{\max}$; their ratio is bounded, and the number of parameter monomials of a given total degree is at most exponential. Enlarge $M,K,d$ to absorb both factors.

For sums and products, use the union of the generator sets and, for a sum of differently shifted supports, the positive difference between their shifts as an additional generator. Coefficient norms are subadditive and submultiplicative, so the same bounds apply. If leading blocks cancel, rebasing at the next nonzero exponent still admits finitely many positive generators: the multi-indices above that exponent form an upward-closed subset of $\N^m$, with finitely many coordinatewise minimal elements. Subtracting the new leading exponent from the weights of those elements supplies finitely many nonnegative shifts, together with the original generators. This also justifies normalizing $F$ by its actual leading monomial in (1). Separate the finite leading polynomial before applying the small-parameter construction; in (2) separate $c+kL$ to obtain the factor $e^cw^k$.
\end{proof}

The theorem supplies sufficient closure conditions, and every truncation is a genuine asymptotic expansion by Lemma~\ref{lem:tail}. They are not necessary conditions for every individual expression: for example $\sqrt{L^2}=-L$ for $L<0$, and polynomials may be applied to an unbounded logarithm. In general, non-integer powers of a nonconstant leading polynomial, its logarithm, exponentials of negative-power arguments, and oscillatory functions of an unbounded logarithm require larger scales (Section~\ref{sec:boundaries}).

\begin{corollary}[Laurent and Puiseux composition]
If $U\in\mathcal E$ has a positive constant leading coefficient and positive valuation, and a selected real branch of $\phi(c+v)$ has a convergent Puiseux expansion $v^{n_0/d}A(v^{1/d})$ for $v\to0^+$, with $d\ge1$ an integer and $A$ analytic at $0$, then $\phi(c+U)\in\mathcal E$. The same holds from below using $\phi(c-v)$.
\end{corollary}
\begin{proof}
Apply the closure theorem to $U^{1/d}$, the analytic function $A$, and $U^{n_0/d}$. Negative $n_0$ includes poles. The side of approach fixes the real branch of a ramified head.
\end{proof}

Variable powers also reduce to the closure theorem through $F^G=\exp(G\log F)$ when $F>0$ and the resulting exponential argument satisfies its hypotheses; this includes $w^w=\exp(w\log w)$. For a nonzero leading block $w^\beta P(L)$, the eventual sign is the sign of $(-1)^{\deg P}[L^{\deg P}]P$, so absolute values reduce to multiplication by that sign whenever it is provable.

\subsection{The precision-tracking evaluator}
\label{sec:evaluator}

The package computes forward expansions by structural recursion over the expression tree, carrying for every subexpression a triple $(T,P,D)$: a jet $T$ and the assertion
\begin{equation}
 \text{subexpression}=T+\OO\bigl(w^P\M(w)^D\bigr),\qquad P\in\R\cup\{+\infty\},
 \label{eq:triple}
\end{equation}
where $P=+\infty$ means that $T$ is exact. A global working order $K_w$ bounds the number of terms of every unit series that is expanded. The rules, all justified by \eqref{eq:O-arithmetic} and Lemma~\ref{lem:tail}, are:
\begin{itemize}[leftmargin=1.6em]
\item constants and $w$ are exact; \wl{Plus}: compute $T_1+T_2$ and truncate below $P=\min(P_1,P_2)$. Combine the input remainder degrees at $P$ with the degree of every explicit block discarded at exponent $P$;
\item \wl{Times}: if $R_i$ denotes the input remainder, the omitted terms include $R_1T_2$, $R_2T_1$ and $R_1R_2$. Their exponent/degree pairs follow from \eqref{eq:O-arithmetic}; use the coarsest pair and truncate $T_1T_2$ below that exponent. Also include the degrees of explicit product blocks discarded exactly at this boundary. A nonempty jet has valuation strictly below its remainder exponent; an empty jet contributes its remainder's exponent and degree when estimating its size;
\item \wl{Power} with an exact real exponent $a$: nonnegative integers by binary exponentiation with precision tracked at each product (exact for exact input). Otherwise the leading block $cw^\beta$ of $T$ must have constant $c>0$ (or $c<0$ with integer $a$); with $U=T/(cw^\beta)-1$, relative precision $P-\beta$ and relative cutoff $c_{\mathrm{rel}}=\min(P-\beta,K_w-a\beta)$, the result is $c^aw^{a\beta}\sum_k\binom akU^k$ with precision $a\beta+c_{\mathrm{rel}}$. For a nonempty $U$, its series-truncation degree is bounded by $\lceil c_{\mathrm{rel}}/\val U\rceil\deg_{\max}U$, combined with the input error degree $D$; an empty $U$ has no computed series tail;
\item \wl{Log}: the leading block must be $cw^\beta$ with constant $c>0$; the result is $\log c+\beta L+\log(1+U)$;
\item \wl{Exp}: $T$ must have no negative exponents and its exponent-$0$ block must be $c+kL$ with exact numeric $k$; the result is $e^cw^k\exp(U)$;
\item an analytic head $\phi$ (\wl{Sin}, \wl{ArcTan}, \wl{BesselJ[0,\,$\cdot$]}, \dots) applied to $T=c+U$ with constant $c$: $\sum_k\phi^{(k)}(c)U^k/k!$ with the Taylor coefficients obtained from \wl{Series[$\phi$[c + t], \{t, 0, N\}]};
\item any other subexpression is expanded by \wl{Series[expr, \{w, 0, n\}, Assumptions -> w > 0]}; the coefficients (which may contain $\log w$) and the irrational-power monomials that \wl{Series} leaves outside its \wl{SeriesData} are parsed as jets, the precision is the order of the \wl{O} term, and the logarithmic degree $D$ is read from the first omitted block obtained by expanding one order further.
\end{itemize}
The whole expression is re-expanded with a larger $K_w$ until the final precision reaches the requested cutoff. The result is reported as the retained blocks plus $\OO_{\sigma_*,k_*}$, where $\sigma_*$ is the exponent of the first omitted block if one lies below the final precision (its complete coefficient having been computed), and otherwise the precision pair $(P,D)$ itself.

\begin{proposition}[Soundness]
\label{prop:soundness}
If every leaf and every Taylor, Laurent, Puiseux or fallback expansion supplied by \wl{Series} satisfies its asserted branch and remainder contract, then the root satisfies \eqref{eq:triple}, and the reported remainder class is valid.
\end{proposition}
\begin{proof}
Induct over the expression tree. For sums and products, use \eqref{eq:O-arithmetic} for propagated errors and add the explicit discarded blocks to the remainder; in particular, their logarithmic degrees at the precision boundary cannot be ignored. For unit series, apply Lemma~\ref{lem:tail} to the finite jet's positive generators and the mean value estimate $\Phi(1+U)-\Phi(1+U_{\mathrm{known}})=\OO(U-U_{\mathrm{known}})$ for its unknown part. Laurent and Puiseux heads are reduced to these same operations on a known signed increment. The final remainder statement follows from the refined tail bound whenever a complete first omitted block has been computed below the precision limit.
\end{proof}

When \wl{Series} supplies a Taylor, Laurent, Puiseux or fallback expansion, its branch and remainder assertions remain inputs to the soundness proposition. Observing the degree of one omitted block alone does not prove that an arbitrary unknown tail satisfies the resulting bound. The structural operations preserve valid assertions, but the package does not prove every analytic contract of a special function or a user-supplied remainder. If cancellation has hidden a leading increment needed by a pole or branch point, the working precision must increase before composition.

\subsection{The second question}

For $F(x)=(1+x+x^{\sqrt2})^{\sqrt2}$ at $x\to+\infty$ the local variable is $w=1/x$ and
\[
 F=w^{-2}\bigl(1+w^{\sqrt2-1}+w^{\sqrt2}\bigr)^{\sqrt2}=w^{-2}\sum_{k\ge0}\binom{\sqrt2}{k}\bigl(w^{\sqrt2-1}+w^{\sqrt2}\bigr)^k ,
\]
an instance of item (1) of Theorem~\ref{thm:closure} with gaps $\sqrt2-1$ and $\sqrt2$. The exponents of $x$ are $2-(\sqrt2-1)k_1-\sqrt2k_2$; ordering the semigroup elements gives \eqref{eq:answer3}, and the first omitted exponent is $7-5\sqrt2$ (multi-index $(5,0)$). The workaround in the question, which peels off one leading term per limit computation, produces the same seven terms; the evaluator supplies them with the remainder $\OO(x^{7-5\sqrt2})$ and can continue to larger finite cutoffs within its resource budget.
